% !TEX program = xelatex
% NOTE: Для корректного Unicode используйте XeLaTeX (или LuaLaTeX).
\documentclass[14pt,a4paper]{extarticle}
\usepackage{fontspec}
\setmainfont{Times New Roman}[
  Ligatures=TeX,
  BoldFont={Times New Roman Bold},
  ItalicFont={Times New Roman Italic},
  BoldItalicFont={Times New Roman Bold Italic}
]
\usepackage[russian]{babel}
\usepackage[left=3cm,right=1cm,top=2cm,bottom=2cm,headheight=1.25cm,footskip=1.25cm]{geometry}
\usepackage{amsmath,amssymb}
\usepackage{array}
\usepackage{booktabs}
\usepackage{longtable}
\usepackage{caption}
\usepackage{enumitem}
\usepackage{setspace}
\usepackage{indentfirst}
\usepackage{titlesec}
\usepackage{tocloft}
\usepackage{hyperref}

\setlength{\parindent}{1.25cm}
\setlength{\parskip}{0pt}
\setstretch{1.5}
\setlist{nosep}
\sloppy
\hyphenpenalty=1000
\exhyphenpenalty=1000
\pagestyle{plain}
\pagenumbering{arabic}
\titleformat{\section}{\normalfont\bfseries\centering}{\thesection}{1em}{\MakeUppercase}
\titleformat{name=\section,numberless}{\normalfont\bfseries\centering}{}{0pt}{\MakeUppercase}
\newcommand{\sectionbreak}{\clearpage}
\titleformat{\subsection}{\normalfont\bfseries}{\thesubsection}{1em}{}
\titlespacing*{\section}{0pt}{0pt}{0pt}
\titlespacing*{\subsection}{1.25cm}{0pt}{0pt}
\titlespacing*{\subsubsection}{1.25cm}{0pt}{0pt}
\titlespacing*{\paragraph}{1.25cm}{0pt}{0pt}
\titlespacing*{\subparagraph}{1.25cm}{0pt}{0pt}
\renewcommand{\cfttoctitlefont}{\hfill\normalfont\bfseries}
\renewcommand{\cftaftertoctitle}{\hfill\mbox{}}
\setlength{\cftsecindent}{1.25cm}
\setlength{\cftsubsecindent}{1.25cm}
\renewcommand{\contentsname}{СОДЕРЖАНИЕ}
\renewcommand{\refname}{Список использованных источников}
\newcolumntype{P}[1]{>{\raggedright\arraybackslash}p{#1}}
\newcolumntype{H}[1]{>{\centering\arraybackslash}p{#1}}
\newlength{\tablecontentwidth}
\DeclareCaptionLabelSeparator{tabledash}{~---~}
\captionsetup[longtable]{position=top,labelsep=tabledash,justification=raggedright,singlelinecheck=false,font=normalsize,labelfont=normalfont,textfont=normalfont,width=\textwidth,skip=0pt}
\setlength{\LTleft}{0pt}
\setlength{\LTright}{0pt plus 1fill}
\setlength{\LTcapwidth}{\textwidth}

\begin{document}
\thispagestyle{empty}

\begin{center}
  {\Large\bfseries Многосторонние конфиденциальные вычисления}\par
\end{center}

\newpage
\tableofcontents

\newpage
\section*{Термины и определения}
\addcontentsline{toc}{section}{Термины и определения}

Многосторонние конфиденциальные вычисления – криптографический подход, позволяющий нескольким участникам совместно вычислять значение функции на приватных входных данных без раскрытия этих данных друг другу.

Протокол – формализованная последовательность действий и сообщений, выполняемых участниками для достижения заданного результата вычисления.

Разделение секрета – метод представления секрета в виде набора долей, при котором отдельная доля не раскрывает исходное значение, а восстановление возможно только при наличии достаточного числа долей.

Зашифрованная схема – представление вычисляемой функции в виде булевой схемы, элементы которой заменены случайными метками и зашифрованными таблицами истинности.

Гомоморфное шифрование – вид шифрования, позволяющий выполнять операции над шифротекстами так, что после расшифрования получается результат соответствующих операций над открытыми данными.

Ослеплённая передача – криптографический протокол, в котором получатель выбирает одно из сообщений отправителя, не раскрывая свой выбор, а отправитель не узнаёт, какое сообщение было получено.

Модель противника – набор предположений о возможностях злоумышленника, включая возможность пассивного наблюдения, отклонения от протокола или преждевременного прекращения вычисления.

Справедливость (fairness) – свойство протокола, при котором результат вычисления получают либо все честные участники, либо ни один из них.

Устойчивость (robustness) – способность протокола сохранять корректность и завершать вычисление при допустимом моделью числе отказавших или действующих злоумышленно участников.

Доверенная среда выполнения (Trusted Execution Environment, TEE) – изолированная аппаратно-программная среда, защищающая исполняемый код и обрабатываемые данные от остальной части системы.

Скрытый доступ к памяти (Oblivious Random Access Memory, ORAM) – метод организации обращений к памяти, скрывающий от наблюдателя адреса и последовательность операций чтения и записи.

Тройки Бивера (Beaver triples) – предварительно сформированные секретно разделённые случайные значения \(a\), \(b\) и \(c\), для которых выполняется равенство \(c=ab\); используются для защищённого умножения секретных величин.

MPC Alliance – международное отраслевое объединение организаций, содействующее распространению, внедрению и стандартизации технологий многосторонних конфиденциальных вычислений.

\newpage
\section*{Перечень сокращений}
\addcontentsline{toc}{section}{Перечень сокращений}

МКВ – многосторонние конфиденциальные вычисления.

МО – машинное обучение.

DP – differential privacy, дифференциальная приватность.

FL – federated learning, федеративное обучение.

GC – garbled circuits, зашифрованные, или запутанные, схемы.

OT – oblivious transfer, ослеплённая передача.

FHE – fully homomorphic encryption, полностью гомоморфное шифрование.

HE – homomorphic encryption, гомоморфное шифрование.

MAC – message authentication code, код аутентификации сообщения.

ORAM – oblivious random access memory, скрытый доступ к памяти.

TEE – trusted execution environment, доверенная среда выполнения.

SGX – software guard extensions, расширения программной защиты Intel.

TDX – trust domain extensions, расширения доверенных доменов Intel.

SEV – secure encrypted virtualization, защищённая зашифрованная виртуализация AMD.

CCA – confidential compute architecture, архитектура конфиденциальных вычислений Arm.

PSI – private set intersection, приватное пересечение множеств.

ReLU – rectified linear unit, функция активации нейронной сети.

zk-SNARK – zero-knowledge succinct non-interactive argument of knowledge, краткий неинтерактивный аргумент знания с нулевым разглашением.

zk-STARK – zero-knowledge scalable transparent argument of knowledge, масштабируемый прозрачный аргумент знания с нулевым разглашением.

\newpage
\section*{Введение}
\addcontentsline{toc}{section}{Введение}

В современном мире актуальна проблема соблюдения конфиденциальности при совместной обработке данных. При совместной работе нескольких организаций задача состоит в получении общего результата без раскрытия приватной информации другим участникам.

Для совместной обработки приватных данных применяются многосторонние конфиденциальные вычисления (МКВ; англ. Multi-\allowbreak Party Computation). Классические модели МКВ позволяют двум или более участникам совместно вычислить значение функции на объединённых входных данных, не раскрывая сами данные и гарантируя корректность результата.

Сферы применения МКВ включает совместную аналитику [17], финансы и здравоохранение. Для задач МО необходимо сочетать защиту данных с требованиями к точности, производительности и масштабируемости вычислений.

Цель работы --- анализ теоретических и криптографических основ МКВ, классификации протоколов и их применений, а также планирование разработки продукта с применением МКВ в области МО.

Для достижения цели поставлены следующие задачи:

\begin{enumerate}

\item Анализ теоретических основ и свойств безопасности МКВ.

\item Исследование криптографических примитивов, используемых в МКВ.

\item Классификация протоколов МКВ.

\item Анализ практических применений МКВ.

\item Выявление ограничений МКВ, условий их практического использования, а также определение направлений развития МКВ.

\item Разработка плана создания и оценки прототипа защищенного МО.

\end{enumerate}

Теоретические основы изложены в разделе~\ref{sec:theoretical-foundations}, криптографические примитивы --- в разделе~\ref{sec:cryptographic-foundations}, классификация --- в разделе~\ref{sec:protocol-classification}, применения --- в разделе~\ref{sec:applications}, ограничения и перспективы --- в разделе~\ref{sec:limitations}, описание проекта защищенного МО --- в разделе~\ref{sec:future-research}; заключение содержит итоги и рекомендации.


\section{Теоретические основы многосторонних конфиденциальных вычислений}
\label{sec:theoretical-foundations}

\subsection{Определение и общая характеристика МКВ}

Формально под многосторонним конфиденциальным вычислением понимается протокол \(\pi\), который принимает от каждого участника \(P_{i}\) его секретный аргумент \(x_{i}\in D\) и совместно с другими участниками вычисляет значение некоторой известной функции \(F : D^{n}\to R\). Протокол должен удовлетворять двум базовым свойствам: \textit{корректность} --- все честные участники получают верный результат \(y=F(x_{1},\ldots,x_{n})\); \textit{конфиденциальность} --- никакая информация о входах честных сторон кроме того, что следует из результата вычислений, не раскрывается другим участникам. Это означает, что протокол обеспечивает информативный эквивалент присутствия доверенной стороны \(T\), которая принимает от всех сторон их входы и возвращает только результат.

Математическая постановка задачи включает описание идеального и реального мира. В \textit{идеальном мире} существует абсолютно доверенная функция \(T\), которой стороны передают свои входы \(x_{i}\). Функция вычисляет \(F(x_{1},\ldots,x_{n})\) и отправляет результат всем сторонам; в этом мире угрозы отсутствуют, а конфиденциальность и корректность гарантированы. В \textit{реальном мире} такой функции нет, и стороны взаимодействуют между собой по протоколу \(\pi\). Цель МКВ --- построить протокол, который делает неразличимыми (с точки зрения вычислительной ограниченности злоумышленника) распределения выходов и представлений (всего, что видит участник или злоумышленник во время выполнения протокола) сторон в идеальном и реальном мирах. В научной литературе это называют \textit{реально-идеальной} парадигмой; она задаёт формальную основу для доказательств безопасности протоколов.

Помимо корректности и конфиденциальности в МКВ выделяют дополнительные свойства: \textit{устойчивость} (robustness), \textit{прекращаемость} и \textit{справедливость} (fairness). Устойчивость означает, что протокол сохраняет корректность вычисления при отказе части участников или при ошибках каналов связи. Справедливость гарантирует, что если одна честная сторона получает результат, то его получают и остальные честные стороны; однако это свойство достижимо только в ограниченных моделях, например при использовании доверенного компонента или экономических стимулов. Поэтому большинство практических протоколов ограничивается гарантиями конфиденциальности и корректности и допускает безопасное досрочное прекращение вычисления при обнаружении атаки.

\subsection{Исторические предпосылки и развитие МКВ}

Истоки МКВ связаны с ментальным покером --- системой криптографических задач 1970-х гг., позволявшей проводить карточные игры без доверенного дилера. Теоретические основы области заложил Эндрю Яо: в 1982 г. он сформулировал задачу миллионеров, а в 1986 г. предложил общий подход к безопасному двустороннему вычислению; термин \textit{искаженная схема} (garbled circuit) закрепился в литературе позднее [14]--[16]. В конце 1980-х гг. появились общие многосторонние протоколы и формальные модели безопасности, показавшие возможность защищённого вычисления произвольных функций [6], [8]. С 2010-х гг. исследования сосредоточены на эффективности и практическом внедрении МКВ, а создание в 2019 г. MPC Alliance отразило переход технологии к отраслевому применению [19].

\subsection{Формальная постановка задачи}

Пусть имеется конечное множество участников \(P_{1},\ldots,P_{n}\). Каждый участник \(P_{i}\) обладает секретным входом \(x_{i}\). Все участники заинтересованы в вычислении результата \(y=F(x_{1},\ldots,x_{n})\), где \(F\) --- детерминированная или вероятностная функция, известная заранее. В идеальной модели с доверенным посредником \(T\) каждый \(P_{i}\) отправляет \(x_{i}\) посреднику, который выдаёт всем результат \(y\). В реальном протоколе отсутствует доверенный посредник, поэтому участники отправляют сообщения друг другу по заранее определённым алгоритмам. Цель протокола \(\pi\) --- обеспечить, чтобы распределения выходных данных и информационных следов сторон в реальном мире были вычислительно неотличимы от таковых в идеальной модели. Для получестного (semi-\allowbreak honest) нарушителя формальное определение описывается через два распределения случайных переменных. Распределение результатов и наблюдений при реальном выполнении протокола обозначается:

\begin{equation}
\mathrm{Real}_{\pi}(\kappa, C; x_{1},\ldots,x_{n})
\end{equation}

где \(\mathrm{Real}_{\pi}\) --- результат выполнения протокола \(\pi\) с параметром безопасности \(\kappa\) и множеством коррумпированных сторон \(C\); распределение включает виды всех коррумпированных сторон и выходы всех участников. В идеальной модели с доверенной функцией и симулятором соответствующее распределение обозначается:

\begin{equation}
\mathrm{Ideal}_{F,\mathrm{Sim}}(\kappa, C; x_{1},\ldots,x_{n})
\end{equation}

где \(\mathrm{Ideal}_{F,\mathrm{Sim}}\) --- распределение в идеальной модели: доверенная функция вычисляет результат, а симулятор \(\mathrm{Sim}\) восстанавливает виды коррумпированных сторон на основе входов и выходов. Протокол безопасно реализует функцию \(F\), если для любого множества коррумпированных сторон эти распределения неразличимы.

Приведём пример: пусть три участника хотят вычислить максимум своих зарплат, не раскрывая их. Они хотят вычислить функцию \(F(x,y,z)=\max(x,y,z)\). В идеальной модели каждый отправляет свой вход доверенному посреднику, который возвращает результат. В реальном мире можно построить протокол на основе разделения секрета: каждый участник делит свой вход на доли и использует операции над зашифрованными битами. Другой вариант --- применить протокол искаженных схем: составить схему сравнения и обменяться зашифрованными метками. Таким образом достигается корректное вычисление \(\max(x,y,z)\) без раскрытия \(x,y,z\) каждому отдельному участнику.

\subsection{Модель вычислений и модели нарушителя}

Модель МКВ задаёт допущения о каналах связи и типах нарушителей. Протоколы МКВ предполагают, что между каждой парой участников имеется аутентифицированный и конфиденциальный канал передачи, что позволяет исключить перехват и модификацию сообщений внешним злоумышленником. Внутренний злоумышленник может нарушать правила протокола, поэтому различают два основных типа:

\begin{itemize}
\item Получестный (semi-\allowbreak honest) нарушитель, или (иногда) честный, но любопытный. Он корректно исполняет инструкции протокола, но пытается по полученным сообщениям восстановить чужие данные. Формально такой нарушитель может объединять представления коррумпированных сторон, но не изменяет ход вычислений. Безопасность протокола обеспечивается наличием симулятора, который в идеальном мире воспроизводит представление нарушителя, что гарантирует отсутствие утечки информации.

\item Злонамеренный (malicious) нарушитель может отклоняться от протокола, отправлять поддельные сообщения, прекращать участие или попытаться изменить вывод честных сторон. Обеспечение безопасности против такого нарушителя требует усиленных механизмов контроля, таких как доказательства с нулевым разглашением, проверка корректности операций и алгоритм "Дели и выбирай" (cut-and-choose). Важно, что многие практические протоколы сначала строятся и анализируются для получестной модели, а затем адаптируются для защиты от злонамеренного нарушителя с использованием протоколов GMW или SPDZ (Shamir-based Protocol for Distributed ZK-proofs) [6].
\end{itemize}

Кроме того, в литературе выделяют модели сетевых задержек (синхронная, полусинхронная, асинхронная), предположения о доле честных участников (например, честное большинство: если \(t\) --- число коррумпированных участников, то \(t<n/2\); такой порог используется, в частности, в некоторых информационно-теоретических протоколах против получестного нарушителя, тогда как для активного нарушителя допустимый порог может быть строже) и возможность наличия дополнительных примитивов (доступ к общему радио-каналу или «чёрной доске») [6], [8]. Эти параметры влияют на возможность достижения критерия справедливости: как показано в ряде работ, для общего числа участников выполнить вычисление с гарантированным выводом без доверенной стороны невозможно, если злоумышленник может прерывать коммуникацию, что известно как теорема Клива [8] --- большинство протоколов реализует только безопасность с возможностью досрочного прекращения.

\subsection{Основные свойства МКВ}

Свойства протокола МКВ можно классифицировать следующим образом:

\begin{enumerate}
\item Конфиденциальность. Никакая информация о входах честных участников, кроме того, что следует из результата, не может быть выведена злоумышленником. Этого достигают за счёт использования разделения секрета, зашифрованных схем, гомоморфного шифрования и других криптографических примитивов.

\item Корректность. Результат, получаемый в ходе протокола, совпадает с результатом идеальной функции \(F\); для получестной модели корректность следует напрямую из определения безопасности, для модели со злонамеренным нарушителем требуются дополнительные проверки.

\item Устойчивость и отказоустойчивость. Протокол способен завершить вычисление при отказе части участников. Многие протоколы на основе разделения секрета допускают восстановление результата при наличии определённого числа честных участников (t-out-of-n).

\item Справедливость. Если хотя бы один честный участник получает вывод, то все честные участники также должны его получить. Достижение справедливости требует специальных механизмов, так как в противном случае злоумышленник может получить результат и прекратить выполнение протокола. В научной литературе гарантированный вывод рассматривается как более сильное свойство, чем справедливость: он влечёт справедливость, однако обратное в общем случае неверно [24]. Его достижение возможно при честном большинстве или с применением дополнительных технологий, таких как защищённые окружения (TEE) или экономические стимулы.
\end{enumerate}

\subsection{Примеры и математические иллюстрации}
\subsubsection{Секретное суммирование}
Пусть \(n\) участников хотят вычислить сумму своих чисел \(s=\sum_{i=1}^{n}x_{i}\). В простой схеме разделения секрета каждый участник выбирает случайные числа \(r_{i}^{(j)}\) для \(j=1\ldots n-1\) так, чтобы сумма этих случайных долей была равна \(x_{i}\). Затем участник передаёт \(r_{i}^{(j)}\) j-му соседу, а оставшуюся долю оставляет себе. Получив все доли, участники локально суммируют полученные значения, и каждый получает долю суммы. Суммируя доли, можно восстановить общий результат, но отдельные входы остаются скрытыми. Такая схема является частным случаем общего механизма разделения секрета, например схемы Шамира, где секрет \(s\) представляет собой значение полинома в точке 0, а доли --- значения этого полинома в точках 1,2, ...

\subsubsection{Схема Шамира}

Рассмотрим схему разделения секрета \((t,n)\) по Шамиру. Дилер выбирает случайный полином степени \(t-1\) над полем \(\mathbb{F}\):

\begin{equation}
f(x)=s+\alpha_{1}x+\alpha_{2}x^{2}+\cdots+\alpha_{t-1}x^{t-1},
\end{equation}

где коэффициенты \(\alpha_{i}\) выбираются случайно. Долю участника \(P_{i}\) определяют как \(s_{i}=f(i)\). Любые \(t\) долей позволяют восстановить секрет с помощью интерполяции Лагранжа, а любой набор из менее чем \(t\) долей не даёт информации о секрете (перфектная конфиденциальность). Эта идея широко используется в протоколах BGW и SPDZ; параметры \(t\) и \(n\) определяют порог устойчивости.

\subsubsection{Искаженные схемы Яо }

Другим механизмом является модель искажённых схем Яо. При её использовании одна сторона («генератор») представляет функцию \(F\) в виде булевой схемы и сопоставляет каждому проводу схемы две случайные метки, соответствующие логическим значениям 0 и 1. Схема преобразуется в «зашифрованную» форму: таблицы заменены на зашифрованные значения. Вторая сторона («вычислитель») с помощью протокола слепой передачи (oblivious transfer) получает метки, соответствующие значениям собственных входных битов, и вычисляет зашифрованную схему, не узнавая открытых значений входов генератора. Полученную выходную метку вычислитель преобразует в результат с помощью таблицы декодирования, предоставленной генератором. Такой подход даёт возможность строить двухсторонние протоколы для произвольных булевых функций, но требует значительных ресурсов памяти и времени.

\subsection{Вывод}

В данном разделе были рассмотрены теоретические основы МКВ: история развития, формальная модель вычислений, модели нарушителя и свойства безопасности; приведены примеры секретного суммирования и схем Яо. Во второй главе будут рассмотрены криптографические основы МКВ: разделение секрета, искаженные схемы, гомоморфное шифрование и протокол слепой передачи.

\section{Криптографические основы многосторонних конфиденциальных вычислений}
\label{sec:cryptographic-foundations}

\subsection{Разделение секрета}

\subsubsection{Определение и свойства}

Разделение секрета --- это криптографическая схема, позволяющая разбить секрет \(S\) на \(n\) частей (долей) \(S_1,\ldots,S_n\) так, что любой набор не менее \(k\) долей позволяет восстановить \(S\), а из меньшего числа долей невозможно получить какую-либо информацию о секрете.  Такой \((k,n)\)-пороговой схеме соответствует два свойства: (1) знание любых \(k\) или более долей делает возможным вычисление секрета; (2) знание не более \(k-1\) долей не уменьшает неопределённость секрета.  Схема считается идеальной и совершенной, если доли имеют ту же длину, что и секрет, и обладают информационно-теоретической стойкостью.

\subsubsection{Схема Шамира}

Наиболее известной реализацией является схема Шамира (Shamir’s secret sharing).  Пусть секрет \(S\) отображается на элемент \(a_0\) конечного поля \(\mathrm{GF}(q)\).  Для порогового значения \(k\) случайно выбирается полином степени \(k-1\)

\begin{equation}
f(x) = a_0 + a_1x + a_2x^2 + \cdots + a_{k-1}x^{k-1},
\end{equation}

где \(a_1,\ldots,a_{k-1}\) --- случайные коэффициенты из \(\mathrm{GF}(q)\).  Для \(i=1,\ldots,n\) вычисляются значения \(f(i)\) и образуют пары \((i, f(i))\); каждому участнику выдаётся одна такая пара.  Любой набор из \(k\) пар позволяет восстановить \(a_0\) с помощью формулы Лагранжа:

\begin{equation}
a_0 = f(0) = \sum_{j=0}^{k-1} y_j \prod_{\substack{m=0\\m \neq j}}^{k-1} \frac{x_m}{x_m - x_j},
\end{equation}

где \((x_j, y_j)\) --- выбранные точки полинома.  Эта формула иллюстрирует, что знание \(k\) различных точек однозначно определяет полином степени \(k-1\) и, следовательно, секрет.

\subsubsection{Пример}

Рассмотрим (3,6)-пороговую схему: секрет \(S=1234\) кодируется полиномом \(f(x)=1234+166x+94x^2\), полученным выбором случайных коэффициентов 166 и 94.  Из полинома вычисляются шесть точек \(f(1)=1494, f(2)=1942,\ldots\), которые выдают шестеро участников.  Любые три точки позволяют восстановить \(f(0)=1234\) через интерполяцию Лагранжа.

\subsubsection{Аддитивное и гомоморфное деление секрета}

Помимо схемы Шамира используется простое аддитивное разделение, при котором секрет разбивается на \(n\) случайных элементов поля так, что их сумма равна \(S\).  Например, для секрета \(s\) выбираются \(n-1\) случайных чисел \(r_1,\ldots,r_{n-1}\) и вычисляется \(r_n = s - \sum_{i=1}^{n-1} r_i\).  В дальнейшем операции над секретами превращаются в поэлементные операции над долями: суммируя доли по модулю \(q\), участники получают долю суммы исходных секретов.  Этот подход лежит в основе протоколов SPDZ и BMR, где для мультипликативных операций используются так называемые triple-векторы Бивера (Beaver triples) [6].  Гомоморфное деление секрета (homomorphic secret sharing) сочетает идеи деления и шифрования и позволяет выполнять локальные операции над долями без дальнейшей коммуникации.

\subsection{Зашифрованные схемы (garbled circuits)}

Зашифрованная схема --- криптографический протокол, позволяющий двум сторонам вычислять любую булеву функцию от своих входных данных без раскрытия этих данных.  Протокол был предложен Яо и требует, чтобы вычисляемая функция была представлена в виде булевой схемы с двумя входами.  Исторический анализ показывает, что термин «garbled circuit» ввели Бивер, Микали и Рогавэй, а идея была описана в работах Голдрайха, Микали и Вигдерсона.

Протокол.  Пусть Алиса (garbler) и Боб (evaluator) хотят вычислить значение функции \(f(x_A, x_B)\).  Протокол состоит из следующих шагов:

\begin{itemize}
\item \textbf{Построение схемы.} Оба участника согласуют булеву схему с 2-входными элементами И/ИЛИ, реализующую функцию \(f\).

\item \textbf{Гарболизация.} Алиса для каждого провода \(w\) в схеме генерирует две случайные \(k\)-битовые метки \(X^w_0, X^w_1\) для логических значений 0 и 1. Затем она проходит по всем элементам и заменяет значения в таблице истинности соответствующими метками. После этого она шифрует каждую строку таблицы двойным симметричным шифрованием с ключами, соответствующими двум входным меткам, и случайно переставляет строки; результат называется garbled table.

\item \textbf{Передача схемы и меток.} Алиса отправляет Бобу зашифрованные таблицы всех элементов и метки, соответствующие своим входам. Метки Боба для каждого бита он получает через протокол ослеплённой передачи (см. § 2.4), что гарантирует, что Алиса не узнает его выбор.

\item \textbf{Оценивание.} Боб, имея garbled circuit и свои метки, проходит по схемным элементам, подбирая для каждого элемента единственную строку таблицы, которую он может расшифровать (двойным шифрованием). Результатом является метка на выходном проводе; Алиса знает соответствие меток фактическим битам и может сообщить результат.
\end{itemize}

В формальном описании схема garbling определяется как случайный алгоритм \(\text{Gb}\), который превращает функцию \(f\) в тройку \((F,e,d)\), где \(e\) --- алгоритм кодирования входов, \(F\) --- зашифрованная функция и \(d\) --- алгоритм декодирования выхода.  При оценивании \(F(e(x))\) получается «garbled output», который после декодирования \(d\) даёт значение \(f(x)\).

Преимущества и ограничения. Зашифрованные схемы подходят для двухсторонних вычислений и обеспечивают однократный обмен данными (небольшое число раундов). Размер схемы и объём передачи пропорциональны количеству элементов И, что увеличивает затраты для сложных функций. Garbled circuits используются для приватного сравнения (миллионеры Яо), проверки знаний и задач с требованиями к низкой задержке [8].

\subsection{Гомоморфное шифрование (homomorphic encryption)}

Гомоморфное шифрование --- это шифр, позволяющий выполнять операции над шифротекстами таким образом, что после расшифрования результат совпадает с результатом той же операции над исходными данными.  Традиционно выделяют частично гомоморфные схемы, поддерживающие одно арифметическое действие (сложение или умножение), и полностью гомоморфные (FHE) схемы, поддерживающие обе операции и, следовательно, реализацию любой булевой функции.  Классическая формулировка утверждает, что гомоморфное шифрование позволяет вычислять над зашифрованными данными без их расшифрования, получая на выходе шифротекст, который при расшифровании даёт результат операции над исходными данными.

Частично гомоморфные шифры.  В простейших схемах выполняется только одно арифметическое действие.  Например, в криптосистеме RSA (без паддинга) шифротекст \(E(m) = m^e \bmod n\) обладает мультипликативным свойством: произведение двух шифротекстов равно шифротексту произведения исходных сообщений.  Схема Эль-Гамаля использует пары \((g^r, m \cdot h^r)\), и перемножение двух шифротекстов приводит к шифрованию произведения сообщений.  В схеме Паилье (Paillier) сообщение шифруется как \(E(m) = g^m r^n \bmod n^2\), и произведение шифротекстов даёт шифр суммы: \(E(m_1)\cdot E(m_2) = E(m_1 + m_2)\).  Эти свойства позволяют строить МКВ-протоколы для простых операций без интерактивности.

Полностью гомоморфное шифрование.  Полностью гомоморфные схемы выполняют как сложение, так и умножение над шифротекстами.  Идея была предложена в 1978 году, но первый практически реализуемый FHE-протокол разработал Гентри в 2009 году, введя метод «bootstrap» для подавления шума.  Схема строится на «частично» гомоморфном шифре, который поддерживает ограниченное число операций, и периодическом гомоморфном декодировании, уменьшающем шум.  Теоретические описания FHE используют абстракцию колец и операции AND/XOR как аналоги умножения/сложения.  Существующие реализации (CKKS, BGV, BFV) основаны на трудности задач «обучения с ошибками» и «приближённого НОД»; они позволяют вычислять целые и вещественные функции над зашифрованными данными, но требуют значительных вычислительных ресурсов и больших размеров ключей.

Преимущества и недостатки. Гомоморфное шифрование позволяет обрабатывать зашифрованный набор данных другого участника без взаимодействия, сокращая число раундов в облачных сервисах и приватной аналитике. Вычислительная сложность FHE и ограничение глубины схем ростом шума обусловливают его сочетание с secret sharing или применение для отдельных операций [2], [5].

\subsection{Протокол ослеплённой передачи (oblivious transfer)}

Определение.  Ослеплённая передача (oblivious transfer, OT) --- это криптографический протокол, в котором отправитель имеет набор сообщений, а получатель получает ровно одно из них, причём отправитель остаётся в неведении о выборе получателя.  Первая версия протокола была предложена М. Рабином, где отправитель посылал сообщение с вероятностью 1/2, и не знал, получил ли его адресат.  Более полезная 1-из-2 OT-схема, разработанная Эвеном, Голдрайхом и Лемпелем, позволяет получателю выбирать один из двух сообщений, не раскрывая свой выбор, и обобщается до 1-из-n OT, где клиент получает один элемент базы данных, ничего не узнавая о других элементах.  Ослеплённая передача является полной для безопасных вычислений: имея реализацию OT, можно безопасно вычислить любую полиномиально вычислимую функцию.

Пример протокола 1-из-2 на основе RSA:
\begin{equation}
\begin{aligned}
v &= (x_b + k^e)\bmod N, \\
m'_i &= \bigl(m_i + (v-x_i)^d\bigr)\bmod N, \quad i\in\{0,1\}, \\
m_b &= (m'_b-k)\bmod N.
\end{aligned}
\end{equation}

Здесь \(N,e,d\) --- параметры RSA, \(x_0,x_1,k\) --- случайные значения, а \(b\in\{0,1\}\) --- выбор Боба, который восстанавливает \(m_b\), не раскрывая Алисе свой выбор и не получая сведений о \(m_{1-b}\).

Роль в МКВ. OT применяется для безопасной передачи входных меток от garbler’а к evaluator’у и в функциональных протоколах, включая Private Set Intersection и oblivious PRF. Возможность реализации произвольных МКВ-протоколов следует из полноты OT.

\subsection{Сравнительный анализ криптографических примитивов}

В таблице~\ref{tab:mpc-primitives} сопоставлены криптографические примитивы МКВ; выбор зависит от модели угроз, числа участников, вычисляемой функции и допустимых затрат.

{\small
\setlength{\tabcolsep}{3pt}
\renewcommand{\arraystretch}{1.2}
\setlength{\tablecontentwidth}{\dimexpr\textwidth-10\tabcolsep-6\arrayrulewidth\relax}
\begin{longtable}{|P{0.16\tablecontentwidth}|P{0.20\tablecontentwidth}|P{0.22\tablecontentwidth}|P{0.22\tablecontentwidth}|P{0.20\tablecontentwidth}|}
\caption{Сравнение криптографических примитивов МКВ}\label{tab:mpc-primitives}\\
\hline
\multicolumn{1}{|H{0.16\tablecontentwidth}|}{\textbf{Примитив}} & \multicolumn{1}{H{0.20\tablecontentwidth}|}{\textbf{Основная идея}} & \multicolumn{1}{H{0.22\tablecontentwidth}|}{\textbf{Преимущество}} & \multicolumn{1}{H{0.22\tablecontentwidth}|}{\textbf{Недостаток}} & \multicolumn{1}{H{0.20\tablecontentwidth}|}{\textbf{Область применения}} \\
\hline
\endfirsthead
\multicolumn{5}{@{}l@{}}{\normalsize Продолжение таблицы~\thetable}\\
\hline
\multicolumn{1}{|H{0.16\tablecontentwidth}|}{\textbf{Примитив}} & \multicolumn{1}{H{0.20\tablecontentwidth}|}{\textbf{Основная идея}} & \multicolumn{1}{H{0.22\tablecontentwidth}|}{\textbf{Преимущество}} & \multicolumn{1}{H{0.22\tablecontentwidth}|}{\textbf{Недостаток}} & \multicolumn{1}{H{0.20\tablecontentwidth}|}{\textbf{Область применения}} \\
\hline
\endhead
Разделение секрета & Секрет кодируется долями; \(k\) из \(n\) долей позволяют восстановить секрет. & Информационно-теоретическая стойкость; лёгкие арифметические операции; эффективная реализация для линейных функций. & Требует много раундов обмена сообщениями для умножения; каждый участник должен быть онлайн. & Протоколы SPDZ/BMR, защищённое машинное обучение, приватная аналитика данных. \\
\hline
Garbled circuits & Функция представляется булевой схемой, её таблицы истинности шифруются случайными метками. & Небольшое число раундов; поддержка произвольных функций; удобство для сравнений. & Большие таблицы и объём передачи; обычно применяется для двух участников; криптографические операции на каждом элементе. & Задача миллионеров Яо, приватное сравнение, доказательства с нулевым разглашением. \\
\hline
Гомоморфное шифрование & Операции выполняются над шифротекстами; после расшифрования результат совпадает с результатом операции над открытыми данными. & Отсутствие интерактивности; одна сторона может выполнить вычисления, не раскрывая входы. & Высокие вычислительные затраты; ограниченная глубина при FHE; большие ключи. & Облачные вычисления, приватные ML-сервисы, вычисления на единичных данных. \\
\hline
Oblivious transfer & Получатель выбирает и получает одно из \(n\) сообщений, не раскрывая выбора; отправитель не узнаёт, какое сообщение выбрано. & Универсальность и полнота для МКВ; удобство для ввода скрытых данных. & Сам по себе не вычисляет функцию; требует криптографических операций, например RSA или протокола Диффи--Хеллмана. & Ввод входов в garbled circuits, PSI (Private Set Intersection), oblivious PRF. \\
\hline
\end{longtable}
}

\subsection{Вывод}

Пороговые схемы разделения секрета обеспечивают информационно-теоретическую защиту и восстановление секрета по долям (\(k\)-кворум); характеристики остальных примитивов сопоставлены в таблице~\ref{tab:mpc-primitives}. В следующем разделе протоколы классифицированы по модели противника, числу участников и криптографическим механизмам.

\section{Классификация протоколов многосторонних конфиденциальных вычислений}
\label{sec:protocol-classification}

\subsection{Классификация по модели противника}
\label{subsec:classification-threats}

Полубезопасная (semi-\allowbreak honest) модель. В этой модели предполагается, что нечестные участники следуют тексту протокола, но пытаются извлечь дополнительную информацию из получаемых сообщений.  Такой противник называется \textit{честный, но любопытный}.  Формальное определение состоит в том, что коррумпированные стороны могут объединять свои виды, но не изменяют сообщения.  Протоколы в полу-честной модели предотвращают непреднамеренную утечку информации и часто служат основой для усиления гарантий; меньшее число криптографических проверок и обменов повышает эффективность.

Активная (malicious) модель. В активной (или злонамеренной) модели противник может отклоняться от протокола, подделывать сообщения, пытаться изменить вывод честных сторон.  Протоколы, обеспечивающие безопасность в этой модели, гарантируют, что при любых действиях противника честные стороны получают корректный результат, а приватность данных сохраняется.  Если при активном нарушителе невозможно вычислить результат (например, нарушитель пытается завершить вычисление преждевременно), честные участники либо обнаруживают факт нарушения и прерывают протокол (\textit{безопасность с абортом}), либо обеспечивается справедливость, когда без результата не остаётся ни одна честная сторона.  Защита от активного противника обычно использует доказательства корректности (zero-knowledge proofs), проверки совпадения входов и копии схем (cut-and-choose), увеличивая издержки.

Скрытная (covert) модель. Между полубезопасной и активной моделями выделяют промежуточный вариант, в котором предполагается, что противник готов нарушить протокол только если его действия останутся незамеченными.  Для таких \textit{скрытных} adversaries протоколы обеспечивают обнаружение нарушения с заданной вероятностью (обычно 75---90 \%).  Модель covert применяется в коммерческих задачах, где репутационные или юридические последствия снижают вероятность атаки, и позволяет повысить эффективность относительно полной активной модели.

\subsection{Классификация по количеству участников}

Двусторонние протоколы (2PC). Двусторонний вариант МКВ предполагает, что только два участника \(P_A\) и \(P_B\) совместно вычисляют функцию \(f(x_A,x_B)\).  Для двух сторон применимы техники, не масштабируемые на множество участников.  Используется схема зашифрованных схем Яо, в которой функция описывается в виде булевой схемы, затем «garbler» отправляет «вычислителю» зашифрованную схему, после чего обе стороны получают результат, ничего более не узнавая.  Яо-протокол обеспечивает безопасность против полудобросовестного противника и характеризуется постоянным числом раундов общения, независимо от сложности функции.  Расширения этого подхода включают активную защиту с cut-and-choose и протокол BMR, который позволяет использовать garbled circuits в многостороннем случае.

Многосторонние протоколы (n-party). Когда вычисление выполняют более двух участников, роли garbler и evaluator исчезают; данные каждой проводки схемы представляются в виде долей, распределённых между сторонами.  В таких протоколах функция изображается как арифметическая схема над конечным полем, а операции сложения и умножения выполняются над секретными долями.  Типичный пример --- протокол GMW, разработанный Голдрайхом, Микали и Вигдерсоном; он использует аддитивное разделение секрета и масштабируется на произвольное число сторон.  GMW допускает активных нарушителей, если большинство участников честное, и вводит доказательства с нулевым разглашением для защиты от обманов.  Другой пример --- протокол BGW, который строится на схеме Шамира и использует арифметические схемы.  Он эффективен для функций, описываемых операциями сложения и умножения, и позволяет устанавливать пороговое значение количества нечестных участников.

Допущение о честном большинстве. В протоколах на основе разделения секрета важную роль играет порог \(t\) числа нечестных сторон.  Для схемы Шамира секрета величина порога, при котором обеспечивается информационно-\allowbreak теоретическая конфиденциальность, зависит от модели противника: при пассивном нарушителе условие \(t< n/2\) достаточно для восстановления секрета, а при активном нарушителе требуется \(t< n/3\).  Аддитивное разделение секрета может выдерживать до \(n-1\) нарушителей в обеих моделях, но требует дополнительных механизмов для проверки корректности (например, MAC-аутентификации в протоколе SPDZ).  Таким образом, классификация по количеству участников тесно связана с предположением об относительном числе честных сторон, что влияет на выбор схемы и гарантии безопасности.

\subsection{Классификация по используемому методу}
\label{subsec:classification-primitives}

В зависимости от используемых криптографических примитивов протоколы МКВ подразделяют на несколько групп. Механизмы и сравнительная характеристика примитивов приведены в разделе~\ref{sec:cryptographic-foundations}. Критерии согласуются с обзором П. Морейры, сопоставляющим протоколы по предположениям безопасности, эффективности, удобству использования и функциональным возможностям [10].

Протоколы на основе разделения секрета используют секретные доли для представления данных. К этой группе относятся BGW, GMW и SPDZ. Гарантии безопасности каждого протокола зависят от принятой модели противника.

К протоколам на основе зашифрованных схем относятся протокол Яо и его многосторонние расширения, например BMR. Отдельную группу составляют протоколы на основе гомоморфного шифрования, в которых данные обрабатываются посредством операций над шифротекстами.

Гибридные протоколы сочетают несколько криптографических примитивов на разных этапах вычисления. Классификация по основному методу не исключает вспомогательных механизмов, в частности ослеплённой передачи (см. раздел~\ref{sec:cryptographic-foundations}).

Специализированные протоколы для приватного пересечения множеств, аукционов и других задач выделяют по назначению, а не по используемому криптографическому методу. Они могут относиться к любой из перечисленных групп. Обзор PSI-протоколов Д. Моралеса, И. Агудо и Х. Лопеса выделяет гомоморфное шифрование, ослеплённую передачу и забывающие псевдослучайные функции и сопоставляет решения по сложности и практической применимости [9]. Классификация по назначению дополняет, но не заменяет критерии примитивов, числа участников и модели противника.

\subsection{Вывод}

Критерии выбора протокола заданы классификацией (см. подразделы~\ref{subsec:classification-threats}--\ref{subsec:classification-primitives}). Двусторонние схемы (garbled circuits, FHE) обеспечивают низкую латентность при сравнении и обработке конфиденциальных данных двух сторон; многосторонние GMW, BGW и SPDZ применяются для распределённых вычислений при честном большинстве. Зависимость коммуникационных и вычислительных затрат от примитивов отражена в таблице~\ref{tab:mpc-primitives}. В следующем разделе рассмотрены применения МКВ в аналитике данных, МО, финансах и медицине.

\section{Применение многосторонних конфиденциальных вычислений в современных системах}
\label{sec:applications}

\subsection{Приватная аналитика данных и аукционы}

В эстонском проекте Sharemind Министерство образования и науки и Налогово-таможенный департамент распределили сведения об обучении и налоговых выплатах между тремя серверами и совместно исследовали связь занятости студентов с завершением обучения. Согласно описанию Cybernetica от 7 июля 2021 г., анализ выполнялся без раскрытия индивидуальных записей и не выявил связи между работой во время учёбы и несвоевременным окончанием вуза, но показал зависимость между уровнем образования и доходом [18].

\subsection{Финансовые приложения и управление ключами}

При распределённом управлении ключом криптовалютного кошелька закрытый ключ разделяется между несколькими участниками, а подпись транзакции формируется совместно без восстановления секрета в одном месте. Такая пороговая схема устраняет единую точку компрометации и позволяет заранее определить число участников, необходимое для подтверждения операции [6].

\subsection{Машинное обучение и анализ больших данных}
\label{subsec:ml-applications}

Общая характеристика. МО часто требует объединения наборов данных; МКВ позволяет обучать модели и выполнять инференс без передачи данных в незащищённом виде. Evans et al. описывают облачное инференс-обслуживание и совместное обучение. В первом случае протокол объединяет скрытый запрос пользователя и параметры модели сервера, возвращая предсказание только пользователю. Во втором несколько организаций обучают модель на объединённых данных без доступа к исходным наборам друг друга, используя секретные доли и часто гибридные методы МКВ и гомоморфного шифрования.

Платформы SecureML, Crypten и Cerebro. Обзор [7] связывает интеграцию МКВ в МО с доступностью библиотек. CRYPTEN исследовательского подразделения Facebook автоматически преобразует модели PyTorch в распределённую МКВ-\allowbreak версию, поддерживает GPU-\allowbreak ускорение и semi-\allowbreak honest модель угроз. Интерфейс, имитирующий API стандартных ML-библиотек, упрощает внедрение. Cerebro реализует доменно-специфический язык описания алгоритмов обучения, оптимизатор и механизмы аудита; компилятор выбирает наиболее эффективные МКВ-протоколы. Механизмы \textit{compute policies} и криптографического аудита контролируют выпуск результатов и корректность вычислений без раскрытия данных.

Обучение с двумя серверами (SecureML).  В классической работе Mohassel \& Zhang (SecureML, 2017) описана схема двухсерверного обучения, при которой данные распределяются между двумя некоррумпированными серверами.  Протокол строит модель линейной или логистической регрессии с использованием стохастического градиентного спуска: в каждой итерации веса \(w_j\) обновляются по правилу

\begin{equation}
w_j := w_j - \alpha \cdot \frac{\partial C_i(w)}{\partial w_j},
\end{equation}

где \(\alpha\) --- шаг обучения, \(C_i\) --- функция потерь для выбранного образца, а \(\partial C_i(w) / \partial w_j\) --- частная производная по параметру \(w_j\).  Для логистической регрессии поверх линейного слоя добавляется сигмоидальная функция активации; нейронные сети строятся как обобщение этих слоёв.  Все арифметические операции выполняются на секретных долях, а сложные функции (например, сигмоидальная активация) аппроксимируются полиномами для сокращения затрат на коммуникацию.  SecureML исключает необходимость доверенного посредника, поскольку данные остаются разделёнными между серверами.

Протоколы глубокого обучения на основе Shamir-\allowbreak деления.  Исследование Zhang et al. (2025) [12] показывает, что для глубокой нейронной сети эффективнее использовать пороговое Shamir-деление.  Авторы предлагают схему с фиксированной точкой для представления вещественных чисел и алгоритм усечения, который ограничивает рост десятичных знаков при умножении; это обеспечивает корректные операции свёртки и функции Softmax, достигая точности более 99 \% при нескольких итерациях.  Предлагаемый протокол оптимизирует умножения с помощью алгоритма Винограда, что снижает число мультипликативных врат более чем на 50 \%.  Основным преимуществом Shamir-базированных схем является более высокая эффективность по сравнению с аддитивным делением, особенно на этапе подготовки и онлайновом умножении; однако поддержание точности требует специальных процедур усечения и поддержки фиксированной запятой.

Проблемы и направления. Ограничения коммуникации, аппроксимации нелинейных функций и глубины схем из-за роста шума рассмотрены в разделе~\ref{sec:limitations}. Разработка библиотек направлена на автоматический подбор протоколов и GPU-ускорение (см. описание платформ выше).

\subsection{Медицина и биоинформатика}

При совместном обучении диагностической модели клиниками данные пациентов остаются локальными, а МКВ защищает вычисление агрегированных обновлений. Обзор [11] относит этот сценарий к основным направлениям конфиденциального медицинского МО; применимость определяется приватностью, вычислительными издержками и масштабируемостью.

\subsection{Блокчейн и децентрализованные приложения}

В протоколе DHSМКВ, предложенном Bao et al. для блокчейн- и облачных систем, многоключевое гомоморфное шифрование позволяет выполнять вычисления над шифротекстами, а направленная расшифровка открывает результат только авторизованной стороне. Подход защищает данные и промежуточные значения без требования взаимного доверия участников [5].

\subsection{Прочие применения: голосование, доказательства с нулевым разглашением и мониторинг}

В электронном голосовании МКВ позволяет конфиденциально суммировать бюллетени; в церемониях формирования параметров доказательств с нулевым разглашением --- распределять генерацию секрета так, чтобы ни один участник не владел им полностью; в сетевом мониторинге --- совместно выявлять совпадения и аномалии без объединения открытых журналов. Во всех сценариях раскрывается только согласованный результат, а исходные данные остаются закрытыми [6].

\subsection{Вывод}

В следующем разделе рассмотрены ограничения внедрения МКВ и способы их преодоления.

\section{Проблемы, ограничения и перспективы развития МКВ}
\label{sec:limitations}

Вычислительные и коммуникационные издержки, скрытый доступ к данным, доверие к среде исполнения, утечки через результаты и организационные барьеры определяют ограничения внедрения МКВ и выбор решений с учётом производительности и модели безопасности.

\subsection{Коммуникационные издержки и гибридные протоколы}
\label{subsec:communication}

Защищённое вычисление требует дополнительных криптографических операций и обменов. При разделении секрета линейные операции часто выполняются локально, а умножение требует взаимодействия. В зашифрованных схемах трафик создаёт передача таблиц нелинейных логических элементов. Рост схемы увеличивает объём обмена, потребление памяти и время выполнения; при высокой сетевой задержке ограничением становится число последовательных раундов [6].

Для снижения издержек гибридные протоколы сочетают разделение секрета, зашифрованные схемы и гомоморфное шифрование: арифметические операции можно выполнять над секретными долями, сравнения --- булевыми схемами, подходящие групповые операции --- над шифротекстами. Представление выбирают по структуре задачи и стоимости переходов между примитивами; их объединение не гарантирует ускорения. Дополнительно применяют пакетную обработку и предварительную генерацию вспомогательных криптографических данных [6, 8].

В МО коммуникационные расходы дополняются трудностями представления дробных чисел, усечения и нелинейных функций. Гибридный подход требует выбора фиксированной точности и оптимизации свёрток и функций активации. В схеме разделения секрета Шамира [12] ускорение оценивается совместно с точностью модели и устойчивостью к переполнению (см. подраздел~\ref{subsec:ml-applications}).

Автоматизированный выбор криптографического представления операций должен учитывать число участников, сетевые задержки и ресурсы, сохраняя заданную модель противника. Сокращение трафика не устраняет преждевременное прекращение протокола: безопасность с прерыванием, справедливость получения результата и устойчивость к отказам остаются самостоятельными требованиями [8].

\subsection{ORAM, масштабируемость и управляемые утечки}

Адреса и последовательность запросов могут раскрывать сведения об операции даже при защите содержимого записей. Для сокрытия шаблонов доступа применяются протоколы ORAM (oblivious random access memory). Дополнительные обращения, обмены и расход памяти усложняют масштабирование МКВ [6].

Накладные расходы можно снизить, допуская управляемую утечку заданных характеристик: размера набора или ограниченных сведений о структуре обращений при сохранении защиты остальных данных. Допустимую утечку необходимо явно включить в модель безопасности и анализировать с учётом дополнительной информации противника.

Выбор между ORAM и управляемыми утечками зависит от чувствительности шаблона обращений. Если его раскрытие недопустимо, издержки снижают организацией данных, пакетной обработкой и оптимизацией скрытого доступа. При допустимом раскрытии упрощение доступа может уменьшить стоимость вычисления. Повторные запросы требуют учёта накопления информации: приемлемость одного раскрытия не гарантирует безопасности последовательности.

Специализированные структуры данных и протоколы согласуют эффективность с заданным профилем утечки. Их оценивают по времени выполнения, памяти и составу раскрываемых сведений; компромисс производительности и утечки не заменяет формального обоснования конфиденциальности.

\subsection{Доверие к среде исполнения и TEE}

Доказательство безопасности протокола не исключает ошибок реализации, компрометации устройств, уязвимостей программно-аппаратного стека обработки долей и побочных каналов утечки. Защита от активного противника не отменяет требований к безопасности устройств и ключевого материала честных участников [6].

Для ограничения доступа к чувствительным данным применяются доверенные среды выполнения (TEE). Они позволяют изолировать выбранные операции от остального программного окружения. В гибридной архитектуре TEE выполняет часть обработки, а МКВ обеспечивает распределённое вычисление. Подход может снижать стоимость операций, но безопасность дополнительно зависит от аппаратной платформы и механизмов её защиты.

Загартдинов и Поляков при сравнении технологий конфиденциальных вычислений подчёркивают необходимость учитывать жизненный цикл защищаемого программного обеспечения и сложности интеграции аппаратных механизмов с существующими средствами защиты [1]. Румянцев и Хаджиева рассматривают защищённую вычислительную среду через взаимосвязь физической и логической составляющих защиты и надёжности вычислений [3].

В настоящей работе из сопоставления этих положений сделан вывод: TEE не устраняет все риски, а её гарантии не следует автоматически приравнивать к криптографическим гарантиям МКВ. Проектирование требует учёта доверия к производителю, проверки состояния кода, управления ключами, обновления компонентов и защиты от побочных каналов. TEE дополняет гибридную архитектуру, но не заменяет МКВ без дополнительных условий.

Направления разработки включают проверяемое развёртывание, аттестацию среды и сочетание аппаратной изоляции с распределением секретов. Архитектура должна определять доступные внутри TEE данные, гарантии при её компрометации, процедуры обновления и восстановления. Применение ускорения требует приемлемой для участников модели доверия.

\subsection{Утечка через результаты и дифференциальная приватность}

МКВ защищают входные данные в процессе вычисления, однако не ограничивают информацию, которую раскрывает разрешённый результат. Точное значение агрегата, предсказание модели или последовательность ответов могут позволить сделать выводы о данных отдельного участника. Такая утечка не обязательно означает нарушение протокола: она может следовать из самой вычисляемой функции [6].

Для ограничения раскрытия через выход МКВ сочетают с дифференциальной приватностью. Результат формируют с контролируемой случайной составляющей, параметры которой выбирают с учётом чувствительности запроса и заданных требований к приватности. Если точный результат нельзя раскрывать ни одной стороне, добавление шума должно быть частью защищённого вычисления до выдачи ответа. Таким образом, МКВ защищают процесс обработки, а дифференциальная приватность ограничивает сведения, доступные из опубликованного результата [6, 13].

В распределённом МО защищённое вычисление с дифференциальной приватностью можно применять к агрегированию обновлений или публикации статистик. Федеративное обучение позволяет оставить исходные данные у их владельцев, но само по себе не устраняет утечки через передаваемые обновления; локальное обучение, защищённое агрегирование и дифференциальная приватность выполняют разные функции, согласуемые в общей архитектуре [13].

Шум может снижать точность модели, а повторные запросы требуют учёта совокупного раскрытия. Компромисс приватности и полезности оценивают через параметры шума и бюджет приватности совместно с точностью, временем выполнения и сетевым трафиком. Необоснованное добавление шума не обеспечивает заявленных гарантий.

\subsection{Правовые и организационные барьеры и стандартизация}
\label{subsec:standardization}

Внедрение МКВ требует согласования обязанностей участников, условий обработки и использования результатов. Криптографическая защита не заменяет правового основания обработки и определения ответственности за ошибки, прекращение вычисления или раскрытие информации. Различия терминологии, моделей безопасности и интерфейсов, а также потребность в специалистах для сопровождения системы на протяжении жизненного цикла создают дополнительные трудности [1, 13].

Стандартизация требует единообразного описания модели противника, допустимого сговора, правил выдачи результата, управляемых утечек и зависимостей от доверенных компонентов. Общие интерфейсы и методики испытаний позволяют сопоставлять совместимость, производительность и безопасность при согласованных предположениях.

Организационные процедуры фиксируют права доступа, подключение участников, обновление ключей и компонентов, действия при отказах и публикацию результатов. Проверка реализации и документации дополняет выбор примитива. Стандартизация не гарантирует правового соответствия: применимость решения оценивается для конкретной задачи и условий обработки.

Направления разработки включают согласование терминологии, профили безопасности типовых приложений и воспроизводимые испытания. Обновление криптографических механизмов, в том числе при изменении модели угроз, предусматривается в жизненном цикле системы. Эти процедуры связывают прототипы с проверяемыми правилами промышленного внедрения; наличие стандарта не заменяет доказательства безопасности.

\subsection{Вывод}

Подходы, рассмотренные в подразделах~\ref{subsec:communication}--\ref{subsec:standardization}, изменяют разные параметры системы: издержки вычисления и скрытого доступа, доверие к среде, раскрытие результатов и сопоставимость требований внедрения.

Ни один подход не устраняет все ограничения; архитектуру оценивают по производительности, точности, допустимому раскрытию и модели доверия. В следующем разделе эти критерии применены к проекту защищённого МО.

\section{Применение МКВ в рамках будущего исследования}
\label{sec:future-research}
\setcounter{table}{3}

\subsection{Постановка задачи и сценарий использования}
\label{subsec:mvp-scenario}

Сценарий предусматривает совместное обучение логистической регрессии для прогнозирования прекращения клиентом использования услуги. Три организации предоставят записи собственных клиентов с одинаковым набором числовых признаков и меткой класса. Исходные данные не должны раскрываться другим организациям и координатору.

Предусматриваются два независимых блока: модель МО и подключаемый МКВ-модуль защищённого исполнения. МКВ-модуль будет исполнять модель на трёх серверах, по одному от каждой организации, без включения логики протокола в модель. Координатор будет запускать задания без доступа к исходным данным.

Модель предполагается обучить на синтетических данных, разделённых между организациями, и проверить инференс на отдельной тестовой выборке. Открытое исполнение той же модели послужит базой сопоставления предсказаний и численных отклонений. Результат инференса будет доступен только владельцу запроса.

Цель MVP --- продемонстрировать обучение и инференс через заменяемый МКВ-модуль без раскрытия входов в рамках semi-honest модели при честном большинстве (см. раздел~\ref{sec:theoretical-foundations}). Демонстрация будет подтверждать работоспособность интеграции, но не заменять доказательство безопасности выбранного протокола.

Проверка максимальной точности, устойчивости к активному противнику, fairness и работы под промышленной нагрузкой в MVP не предусматривается. Выбор логистической регрессии и готовых библиотечных операций будет обусловлен простотой реализации и воспроизводимостью, а не максимальной эффективностью или защищённостью. Прототип не предполагается использовать как промышленное решение.

\subsection{Архитектура продукта}
\label{subsec:product-architecture}

Продукт предполагается разделить на блок МО, задающий модель и порядок обучения, и МКВ-модуль защищённого исполнения. Подготовка данных и координация заданий будут обслуживающими функциями, а не самостоятельными вычислительными блоками. Разделение позволит проверять модель в открытом режиме и подключать защищённое исполнение через согласованный интерфейс. Приоритетом MVP будет простота интеграции, а не универсальная поддержка моделей.

Блок МО будет содержать описание логистической регрессии, порядок признаков и правила обновления весов. Вход обучения --- матрица числовых признаков и вектор бинарных меток; вход инференса --- запись того же формата без метки. Выходом станет оценка вероятности прекращения использования услуги. Владелец запроса преобразует её в класс по установленному порогу после получения результата. Для упрощения MVP это исключит отдельное защищённое сравнение при выдаче ответа.

МКВ-модуль будет выполнять операции над защищёнными значениями на трёх серверах. Описание модели останется публичным; входы, метки, веса и промежуточные результаты будут представлены долями. Предварительное вычисление открытого предсказания с последующим разделением результата не допускается. Криптографические операции будут отделены от описания модели; их основания и сопоставление приведены в разделе~\ref{sec:cryptographic-foundations} и таблице~\ref{tab:mpc-primitives}.

Интерфейс блоков будет включать инициализацию параметров, обучение, инференс и сохранение состояния. Публичная конфигурация задаст версию модели, размерности, параметры квантования, число итераций и получателя результата. Приватные входы и состояние будут передаваться через ссылки на защищённые объекты, без открытых массивов. Модуль вернёт ссылку на обновлённое состояние, статус и результат для разрешённого получателя. Блок МО будет задавать операции, но не управлять обменом и восстановлением значений.

Скрытие данных от модели означает отсутствие открытых значений у кода, описывающего вычисления, а не защиту от модели как самостоятельного противника. Этот код будет обращаться к защищённым объектам через МКВ-модуль. Другие организации не получат исходные записи; каждому серверу будут переданы предусмотренные протоколом доли. Границы доступа приведены в таблице~\ref{tab:mvp-architecture}.

{\small
\setlength{\tabcolsep}{3pt}
\renewcommand{\arraystretch}{1.2}
\setlength{\tablecontentwidth}{\dimexpr\textwidth-6\tabcolsep-4\arrayrulewidth\relax}
\begin{longtable}{|P{0.22\tablecontentwidth}|P{0.39\tablecontentwidth}|P{0.39\tablecontentwidth}|}
\caption{Компоненты продукта и границы раскрытия данных}\label{tab:mvp-architecture}\\
\hline
\multicolumn{1}{|H{0.22\tablecontentwidth}|}{\textbf{Компонент}} & \multicolumn{1}{H{0.39\tablecontentwidth}|}{\textbf{Функция}} & \multicolumn{1}{H{0.39\tablecontentwidth}|}{\textbf{Доступные данные}} \\
\hline
\endfirsthead
\multicolumn{3}{@{}l@{}}{\normalsize Продолжение таблицы~\thetable}\\
\hline
\multicolumn{1}{|H{0.22\tablecontentwidth}|}{\textbf{Компонент}} & \multicolumn{1}{H{0.39\tablecontentwidth}|}{\textbf{Функция}} & \multicolumn{1}{H{0.39\tablecontentwidth}|}{\textbf{Доступные данные}} \\
\hline
\endhead
Клиент организации & Проверка формата, подготовка и квантование входа; передача долей серверам. & Собственные записи и метки; публичная конфигурация задания. \\
\hline
Блок МО & Описание модели, последовательности операций и правил обучения. & Публичное описание; ссылки на защищённые значения, без их раскрытия. \\
\hline
МКВ-модуль: три сервера & Защищённое обучение и инференс; хранение состояния; выдача результата. & Локальные доли входов, весов и промежуточных значений; сообщения протокола. \\
\hline
Координатор & Запуск задания и сбор статусов. & Публичные настройки и служебные сведения; без записей, долей и весов. \\
\hline
Владелец запроса & Восстановление результата инференса и применение порога классификации. & Собственный запрос и его предсказание; без чужих входов и весов. \\
\hline
\end{longtable}
}

Поток данных: клиент $\rightarrow$ локальная подготовка и квантование $\rightarrow$ МКВ-модуль на трёх серверах $\rightarrow$ результат владельцу запроса. Блок МО передаст описание вычислений, координатор --- управляющее задание. Порядок признаков, масштабирование и границы значений будут согласованы до запуска. Заранее заданные параметры нормализации позволят упростить MVP, исключив отдельный протокол расчёта статистик объединённой выборки.

Веса после обучения будут сохраняться распределённо: каждый сервер будет хранить собственную долю и идентификатор конфигурации. Последующий инференс будет использовать это состояние без восстановления полного набора весов у координатора. Открытая контрольная модель будет храниться отдельно и использоваться только в эксперименте с синтетическими данными. Журналы будут содержать статусы, размерности и время выполнения, но не записи, доли, градиенты или предсказания; скрытие служебных сведений в MVP не предусматривается.

Заменяемость модуля потребует совместимости интерфейса, числового представления и правил раскрытия. Для упрощения MVP предусматривается адаптер с фиксированными операциями логистической регрессии, а не преобразователь произвольных сетей. Замена реализации потребует проверки численной совместимости и модели угроз, но не изменения описания модели. Оптимизация обмена и расширение вычислений будут рассматриваться после демонстрации работоспособности с учётом раздела~\ref{sec:limitations}.

\subsection{Модель угроз и границы MVP}
\label{subsec:mvp-threat-model}

Модель угроз прототипа будет конкретизировать постановку из пп.~1.3--1.4 применительно к архитектуре подраздела~\ref{subsec:product-architecture}. Предусматриваются три организации и три вычислительных сервера, по одному от каждой организации. Координатор будет управлять заданиями без доступа к приватным входам и долям. Допущения приведены в таблице~\ref{tab:mvp-threat-model}.

{\small
\setlength{\tabcolsep}{3pt}
\renewcommand{\arraystretch}{1.2}
\setlength{\tablecontentwidth}{\dimexpr\textwidth-4\tabcolsep-3\arrayrulewidth\relax}
\begin{longtable}{|P{0.27\tablecontentwidth}|P{0.73\tablecontentwidth}|}
\caption{Допущения модели угроз MVP}\label{tab:mvp-threat-model}\\
\hline
\multicolumn{1}{|H{0.27\tablecontentwidth}|}{\textbf{Параметр}} & \multicolumn{1}{H{0.73\tablecontentwidth}|}{\textbf{Принятое допущение}} \\
\hline
\endfirsthead
\multicolumn{2}{@{}l@{}}{\normalsize Продолжение таблицы~\thetable}\\
\hline
\multicolumn{1}{|H{0.27\tablecontentwidth}|}{\textbf{Параметр}} & \multicolumn{1}{H{0.73\tablecontentwidth}|}{\textbf{Принятое допущение}} \\
\hline
\endhead
Противник & semi-honest: наблюдает состояние своего сервера и сообщения, но не изменяет предписанные операции. \\
\hline
Допустимый сговор & Не более одной организации (её клиент и сервер) за весь цикл работы. Допускается сговор с координатором. \\
\hline
Честное большинство & Не менее двух серверов честны. Административные области независимы; общий доступ к состояниям серверов исключён. \\
\hline
Каналы и среда & Аутентифицированные конфиденциальные каналы; корректная реализация и случайность. Компрометация среды и побочные каналы исключены из модели. \\
\hline
Допустимое раскрытие & Конфигурация, размеры выборок, время и статусы. Предсказание раскрывается владельцу запроса; веса и промежуточные значения не восстанавливаются. \\
\hline
\end{longtable}
}

Защите будут подлежать чужие записи, метки и промежуточные значения; собственные входы противнику уже известны. Допустимость сговора является требованием к протоколу, а не следствием наличия трёх серверов. Зависимости от вспомогательных компонентов подготовки вычислений будут учитываться при выборе платформы. Три процесса на одном компьютере позволят проверить интеграцию, но не независимость административных областей.

Защита от активного противника, проверка истинности данных, предотвращение подмены долей и намеренного прекращения участия не входят в MVP. Выбор semi-honest модели позволит упростить реализацию, исключив проверки активной безопасности, а не обеспечить максимальную защищённость. При отказе сервера задание завершится с ошибкой; устойчивость к отказам, гарантированная выдача результата и fairness не заявляются.

MVP будет ограничен логистической регрессией и синтетическими данными, без поддержки сетей произвольной глубины, промышленной нагрузки и автоматического масштабирования. Вывод сведений об обучающей выборке из разрешённых предсказаний, в том числе при повторных запросах, остаётся вне заявленной защиты. DP и скрытие служебных сведений не предусматриваются; эти ограничения рассмотрены в разделе~\ref{sec:limitations}. Гарантии за пределами принятых допущений не заявляются.

\subsection{Выбор протокола и платформы}
\label{subsec:mvp-protocol-platform}

Согласно таблице~\ref{tab:mpc-primitives}, для арифметической части модели предполагается использовать разделение секрета. В качестве базового сочетания выбираются TFEncrypted и semi-honest вариант ABY3 для трёх серверов \cite{20,21}. Выбор будет обусловлен доступностью библиотечных операций и соответствием допущениям подраздела~\ref{subsec:mvp-threat-model}, а не максимальной производительностью или защищённостью. Обзор применений платформ приведён в подразделе~\ref{subsec:ml-applications}; их сопоставление для данного MVP представлено в таблице~\ref{tab:mvp-platforms}.

{\small
\setlength{\tabcolsep}{3pt}
\renewcommand{\arraystretch}{1.2}
\setlength{\tablecontentwidth}{\dimexpr\textwidth-6\tabcolsep-4\arrayrulewidth\relax}
\begin{longtable}{|P{0.19\tablecontentwidth}|P{0.39\tablecontentwidth}|P{0.42\tablecontentwidth}|}
\caption{Сопоставление платформ для реализации MVP}\label{tab:mvp-platforms}\\
\hline
\multicolumn{1}{|H{0.19\tablecontentwidth}|}{\textbf{Платформа}} & \multicolumn{1}{H{0.39\tablecontentwidth}|}{\textbf{Основание применения}} & \multicolumn{1}{H{0.42\tablecontentwidth}|}{\textbf{Решение для MVP}} \\
\hline
\endfirsthead
\multicolumn{3}{@{}l@{}}{\normalsize Продолжение таблицы~\thetable}\\
\hline
\multicolumn{1}{|H{0.19\tablecontentwidth}|}{\textbf{Платформа}} & \multicolumn{1}{H{0.39\tablecontentwidth}|}{\textbf{Основание применения}} & \multicolumn{1}{H{0.42\tablecontentwidth}|}{\textbf{Решение для MVP}} \\
\hline
\endhead
CrypTen \cite{22} & Интерфейс на основе PyTorch; арифметические и бинарные доли. & Не выбран: режим с доверенным генератором троек добавит допущение. Репозиторий архивирован. \\
\hline
TFEncrypted \cite{21} & Интерфейс TensorFlow; ABY3 на трёх серверах; операции обучения и инференса. & Выбран: готовые операции соответствуют составу участников. Потребуется фиксация совместимых версий. \\
\hline
SecureML \cite{23} & Исследовательская реализация МО на двух не вступающих в сговор серверах. & Не выбран: перенос в заданную трёхсерверную архитектуру потребует дополнительных решений. \\
\hline
\end{longtable}
}

Для линейных операций будет использоваться реплицированное аддитивное разделение в ABY3, без самостоятельной реализации схемы Шамира \cite{20,21}. Сложение, умножение матриц и усечение фиксированной точки будут выполняться библиотекой. Применение готового варианта сократит число собственных криптографических процедур; переход к иной схеме не требуется для демонстрации обучения и инференса. Представление чисел и допустимые отклонения будут зафиксированы в требованиях к прототипу.

Для сигмоиды выбирается библиотечная аппроксимация с тремя линейными участками, предусмотренная реализацией ABY3 \cite{21}. Это проще разработки собственного алгоритма нелинейных вычислений. Выбор участка будет выполняться защищёнными сравнениями над бинарными долями и библиотечными преобразованиями представлений, без раскрытия результатов сравнений. Интеграция отдельной подсистемы GC не предусматривается; используемый библиотекой OT не будет реализовываться заново. Такая конфигурация сохраняет нелинейную обработку внутри МКВ-модуля, а не переносит её на открытые значения.

Для обучения предполагается фиксированное число шагов обновления весов без защищённого критерия остановки. Открытая контрольная реализация будет применять ту же аппроксимацию; её влияние на качество будет дополнительно оцениваться относительно исходной сигмоиды. Это позволит отделить отклонения квантования и защищённого исполнения от изменения самой модели. Пороговая классификация останется локальной операцией владельца результата согласно подразделу~\ref{subsec:product-architecture}.

FHE, самостоятельная интеграция HE и GC, а также протоколы защиты от активного противника не выбираются для MVP. Они потребуют дополнительных механизмов и настроек, избыточных относительно проверки работоспособности. Это решение следует ограничениям таблицы~\ref{tab:mpc-primitives} и принятой модели угроз, а не утверждает непригодность этих подходов для других задач.

Подготовка случайности будет выполняться средствами выбранного трёхстороннего протокола, без отдельного доверенного генератора троек и без передачи этой роли координатору \cite{20,21}. Для воспроизводимости будет зафиксирована версия TFEncrypted с проверенными операциями; наличие криптографической случайности, раздельное размещение серверов и защищённые каналы будут проверяться отдельно. TFEncrypted позиционируется как экспериментальное программное обеспечение \cite{21}; его выбор не заменяет проверки реализации. Адаптер из подраздела~\ref{subsec:product-architecture} позволит заменить платформу без изменения описания модели.

\subsection{Функциональные и нефункциональные требования}
\label{subsec:mvp-requirements}

Требования будут определять проверяемый объём прототипа в рамках подразделов~\ref{subsec:mvp-scenario}--\ref{subsec:mvp-protocol-platform}. Их выполнение должно подтвердить обучение и инференс через отдельный МКВ-модуль, а не достижение промышленной производительности. Функциональные требования приведены в таблице~\ref{tab:mvp-functional-requirements}; форматы и параметры будут фиксироваться общей конфигурацией задания.

{\small
\setlength{\tabcolsep}{3pt}
\renewcommand{\arraystretch}{1.2}
\setlength{\tablecontentwidth}{\dimexpr\textwidth-4\tabcolsep-3\arrayrulewidth\relax}
\begin{longtable}{|P{0.12\tablecontentwidth}|P{0.88\tablecontentwidth}|}
\caption{Функциональные требования к MVP}\label{tab:mvp-functional-requirements}\\
\hline
\multicolumn{1}{|H{0.12\tablecontentwidth}|}{\textbf{Код}} & \multicolumn{1}{H{0.88\tablecontentwidth}|}{\textbf{Требование}} \\
\hline
\endfirsthead
\multicolumn{2}{@{}l@{}}{\normalsize Продолжение таблицы~\thetable}\\
\hline
\multicolumn{1}{|H{0.12\tablecontentwidth}|}{\textbf{Код}} & \multicolumn{1}{H{0.88\tablecontentwidth}|}{\textbf{Требование}} \\
\hline
\endhead
Ф1 & Система должна принимать локальные файлы CSV с согласованными признаками и, при обучении, метками; проверять размерности, диапазоны и отсутствие пропусков до передачи данных. Несоответствующий вход должен отклоняться. \\
\hline
Ф2 & Система должна выполнять квантование и разделение входов на стороне владельца, передавать доли трём серверам и не передавать исходные записи или доли координатору. \\
\hline
Ф3 & Система должна обучать логистическую регрессию через МКВ-модуль с заданным числом обновлений весов и библиотечной аппроксимацией сигмоиды, без включения криптографической логики в описание модели. \\
\hline
Ф4 & Система должна сохранять и загружать состояние модели как локальные доли на серверах с идентификатором конфигурации, без восстановления полного набора весов. \\
\hline
Ф5 & Система должна выполнять защищённый инференс и восстанавливать оценку вероятности только у владельца запроса. Преобразование в класс должно выполняться локально после выдачи результата. \\
\hline
Ф6 & Система должна обеспечивать открытое контрольное исполнение на синтетических данных и сопоставление с разрешёнными результатами защищённого режима; формировать отчёт о качестве, времени и трафике. \\
\hline
\end{longtable}
}

Нефункциональные требования представлены в таблице~\ref{tab:mvp-nonfunctional-requirements}. Значения являются проектными целями, а не результатами измерений. Для упрощения MVP будет использоваться одна конфигурация модели и числового представления, без автоматического подбора параметров.

{\small
\setlength{\tabcolsep}{3pt}
\renewcommand{\arraystretch}{1.2}
\setlength{\tablecontentwidth}{\dimexpr\textwidth-4\tabcolsep-3\arrayrulewidth\relax}
\begin{longtable}{|P{0.25\tablecontentwidth}|P{0.75\tablecontentwidth}|}
\caption{Нефункциональные требования и целевые параметры MVP}\label{tab:mvp-nonfunctional-requirements}\\
\hline
\multicolumn{1}{|H{0.25\tablecontentwidth}|}{\textbf{Показатель}} & \multicolumn{1}{H{0.75\tablecontentwidth}|}{\textbf{Цель и условия проверки}} \\
\hline
\endfirsthead
\multicolumn{2}{@{}l@{}}{\normalsize Продолжение таблицы~\thetable}\\
\hline
\multicolumn{1}{|H{0.25\tablecontentwidth}|}{\textbf{Показатель}} & \multicolumn{1}{H{0.75\tablecontentwidth}|}{\textbf{Цель и условия проверки}} \\
\hline
\endhead
Состав системы & Три организации, три вычислительных сервера и один координатор; модель угроз --- согласно подразделу~\ref{subsec:mvp-threat-model}. \\
\hline
Контрольные данные & 900 обучающих записей, по 300 у каждой организации; 300 отдельных тестовых записей. 20 числовых признаков, бинарная метка; признаки в диапазоне $[-1;1]$. \\
\hline
Числовой формат & Фиксированная точка: 64-битное представление, 16 дробных битов, масштаб $2^{16}$, в соответствии с доступной конфигурацией библиотеки \cite{21}. Усечение --- средствами МКВ-модуля. \\
\hline
Режим обучения & Нулевые начальные веса; шаг обучения $1/16$; 150 обновлений с размером пакета 30 записей, по 10 от каждой организации. Порог классификации --- $0{,}5$. \\
\hline
Численная точность и качество & Снижение доли правильных классификаций --- не более 2 процентных пунктов относительно открытой модели с той же аппроксимацией. Среднее абсолютное отклонение вероятностей --- не более $0{,}005$. \\
\hline
Сетевой обмен & Ориентировочный бюджет --- не более 10~МиБ в среднем на один шаг обучения для указанного пакета. Учитываются исходящие сообщения трёх серверов, включая подготовку, выполняемую внутри шага. \\
\hline
Испытательная среда & Три узла, каждый не менее двух процессорных ядер и 8~ГиБ оперативной памяти; сеть не менее 100~Мбит/с. Время обучения и инференса измеряется без обязательного предельного значения. \\
\hline
Воспроизводимость & Пять повторных защищённых запусков с одинаковыми данными, порядком пакетов и настройками. Фиксируются версии зависимостей и параметры генерации синтетических данных. \\
\hline
\end{longtable}
}

Точность будет измеряться на одной тестовой выборке при одинаковом пороге классификации. Предел в 2 процентных пункта означает абсолютную разность долей правильных ответов, а не относительное снижение на 2~\%. Оба критерия качества должны выполняться в каждом из пяти запусков. Сопоставление с открытой моделью, использующей исходную сигмоиду, будет приводиться отдельно, чтобы учитывать эффект аппроксимации согласно подразделу~\ref{subsec:mvp-protocol-platform}.

Объём обмена будет суммироваться по отправленным сообщениям без повторного учёта на стороне получателей. Первоначальная загрузка входов и настройка соединений будут измеряться отдельно. Превышение бюджета трафика должно отражаться в отчёте, но не требовать перехода к более сложному протоколу ради оптимизации MVP. Параметры среды будут фиксироваться для сопоставимости измерений.

Достаточность разрядности будет проверяться до защищённого запуска для заданных диапазонов и числа шагов, без раскрытия промежуточных значений. Журналы не должны содержать входы, доли, веса, градиенты или предсказания. Генератор синтетических данных и криптографическая случайность будут инициализироваться независимо; повторное использование секретной случайности между запусками недопустимо.

\subsection{План реализации и критерии готовности MVP}
\label{subsec:mvp-implementation-plan}

Этапы реализации приведены в таблице~\ref{tab:mvp-implementation-stages}. Для упрощения MVP будут использоваться готовые операции с интерфейсом подраздела~\ref{subsec:product-architecture}, без собственных криптографических процедур. Данные и настройки задаются таблицей~\ref{tab:mvp-nonfunctional-requirements}.

Бинарные метки синтетических данных будут задаваться порогом линейной комбинации признаков; выборки будут генерироваться раздельно. Такая зависимость упростит проверку обучения, но не будет воспроизводить статистику реальных клиентов.

{\small
\setlength{\tabcolsep}{3pt}
\renewcommand{\arraystretch}{1.2}
\setlength{\tablecontentwidth}{\dimexpr\textwidth-6\tabcolsep-4\arrayrulewidth\relax}
\begin{longtable}{|P{0.22\tablecontentwidth}|P{0.40\tablecontentwidth}|P{0.38\tablecontentwidth}|}
\caption{Этапы реализации MVP и ожидаемые результаты}\label{tab:mvp-implementation-stages}\\
\hline
\multicolumn{1}{|H{0.22\tablecontentwidth}|}{\textbf{Этап}} & \multicolumn{1}{H{0.40\tablecontentwidth}|}{\textbf{Содержание работы}} & \multicolumn{1}{H{0.38\tablecontentwidth}|}{\textbf{Ожидаемый результат}} \\
\hline
\endfirsthead
\multicolumn{3}{@{}l@{}}{\normalsize Продолжение таблицы~\thetable}\\
\hline
\multicolumn{1}{|H{0.22\tablecontentwidth}|}{\textbf{Этап}} & \multicolumn{1}{H{0.40\tablecontentwidth}|}{\textbf{Содержание работы}} & \multicolumn{1}{H{0.38\tablecontentwidth}|}{\textbf{Ожидаемый результат}} \\
\hline
\endhead
Подготовка данных & Генерация выборок, распределение 900 обучающих записей между тремя владельцами; запуск открытой модели. & Три файла обучения, 300 тестовых записей и контрольные предсказания. \\
\hline
Квантование & Локальное преобразование входов в фиксированную точку; проверка диапазонов и граничных случаев. & Согласованный формат; открытые проверки кодирования и арифметики пройдены. \\
\hline
Разделение на доли & Подключение библиотечного ввода и передача долей серверам без участия координатора в обработке значений. & На серверах размещены локальные доли; исходные записи остаются у владельцев. \\
\hline
Вычисление & 150 обновлений весов; сохранение состояния; инференс после перезапуска серверных процессов. & Загружаемые доли весов и завершённое вычисление 300 тестовых ответов. \\
\hline
Восстановление результата & Выдача вероятностей владельцу запроса; локальная классификация и сопоставление с открытым исполнением. & Разрешённые предсказания и отчёт о качестве, времени и обмене. \\
\hline
\end{longtable}
}

MVP будет считаться готовым при выполнении всех условий таблицы~\ref{tab:mvp-readiness}. Проверки будут проводиться в зафиксированном окружении на трёх раздельных вычислительных узлах. Локальный запуск трёх процессов будет использоваться только при отладке, согласно ограничениям подраздела~\ref{subsec:mvp-threat-model}.

{\small
\setlength{\tabcolsep}{3pt}
\renewcommand{\arraystretch}{1.2}
\setlength{\tablecontentwidth}{\dimexpr\textwidth-4\tabcolsep-3\arrayrulewidth\relax}
\begin{longtable}{|P{0.25\tablecontentwidth}|P{0.75\tablecontentwidth}|}
\caption{Критерии готовности MVP}\label{tab:mvp-readiness}\\
\hline
\multicolumn{1}{|H{0.25\tablecontentwidth}|}{\textbf{Проверка}} & \multicolumn{1}{H{0.75\tablecontentwidth}|}{\textbf{Условие приёмки}} \\
\hline
\endfirsthead
\multicolumn{2}{@{}l@{}}{\normalsize Продолжение таблицы~\thetable}\\
\hline
\multicolumn{1}{|H{0.25\tablecontentwidth}|}{\textbf{Проверка}} & \multicolumn{1}{H{0.75\tablecontentwidth}|}{\textbf{Условие приёмки}} \\
\hline
\endhead
Запуск и входы & Документированный сценарий запускает обучение и инференс на заданных 900 и 300 записях с 20 признаками. Неверная размерность, пропуски и выход за диапазон отклоняются до передачи долей. \\
\hline
Обучение и состояние & Выполнены 150 обновлений от нулевых весов. После сохранения и перезапуска инференс использует загруженные доли без повторного обучения. Точность выше постоянного прогноза преобладающего в обучающей выборке класса. \\
\hline
Численная совместимость & В каждом из пяти запусков на 300 тестовых записях просадка относительно открытой модели с той же аппроксимацией не превышает 2 процентных пунктов; среднее абсолютное отклонение вероятностей не превышает $0{,}005$. \\
\hline
Границы раскрытия & Проверка интерфейсов и трасс подтверждает отсутствие восстановления входов, весов и градиентов. Предсказания выдаются только владельцу запроса; журналы не содержат запрещённых значений из подраздела~\ref{subsec:mvp-requirements}. \\
\hline
Отчёт и воспроизведение & Для пяти запусков сохранены конфигурации, версии и агрегированные метрики. Трафик шага отделён от загрузки входов; сопоставление с бюджетом 10~МиБ отражено в отчёте. \\
\hline
\end{longtable}
}

Проверка интерфейсов и трасс не заменяет доказательство безопасности. Готовность определяется в рамках подраздела~\ref{subsec:mvp-threat-model}; превышение бюджета обмена отражается в отчёте без усложнения протокола.

\subsection{Риски, ограничения и направления развития после MVP}
\label{subsec:mvp-risks-roadmap}

Ограничения прототипа будут определяться подразделом~\ref{subsec:mvp-threat-model} и выводами раздела~\ref{sec:limitations}. Направления таблицы~\ref{tab:mvp-risks-roadmap} не входят в критерии готовности MVP: их реализация увеличит сложность относительно цели демонстрации. Риски несовместимости зависимостей и численного представления будут контролироваться фиксацией версий и проверками подразделов~\ref{subsec:mvp-requirements}--\ref{subsec:mvp-implementation-plan}, а не расширением набора платформ.

{\small
\setlength{\tabcolsep}{3pt}
\renewcommand{\arraystretch}{1.2}
\setlength{\tablecontentwidth}{\dimexpr\textwidth-6\tabcolsep-4\arrayrulewidth\relax}
\begin{longtable}{|P{0.23\tablecontentwidth}|P{0.34\tablecontentwidth}|P{0.43\tablecontentwidth}|}
\caption{Ограничения MVP и направления последующего развития}\label{tab:mvp-risks-roadmap}\\
\hline
\multicolumn{1}{|H{0.23\tablecontentwidth}|}{\textbf{Область риска}} & \multicolumn{1}{H{0.34\tablecontentwidth}|}{\textbf{Граница MVP}} & \multicolumn{1}{H{0.43\tablecontentwidth}|}{\textbf{Развитие после MVP}} \\
\hline
\endfirsthead
\multicolumn{3}{@{}l@{}}{\normalsize Продолжение таблицы~\thetable}\\
\hline
\multicolumn{1}{|H{0.23\tablecontentwidth}|}{\textbf{Область риска}} & \multicolumn{1}{H{0.34\tablecontentwidth}|}{\textbf{Граница MVP}} & \multicolumn{1}{H{0.43\tablecontentwidth}|}{\textbf{Развитие после MVP}} \\
\hline
\endhead
Точность и издержки & Синтетические данные и одна модель; масштабирование и промышленная нагрузка не проверяются. & Расширить испытания; проверить разрядность новых операций; затем оптимизировать пакетный обмен согласно подразделу~\ref{subsec:communication}. \\
\hline
Модель противника & Активные нарушения и сговор двух организаций вне модели; fairness не обеспечивается. & Выбрать иной профиль безопасности и заменить реализацию МКВ-модуля. Fairness и устойчивость проверять как отдельные свойства. \\
\hline
Раскрытие через ответы & DP и защита от вывода сведений по повторным запросам не реализуются. & Исследовать DP внутри МКВ-модуля до выдачи ответа; учитывать бюджет приватности и изменение качества согласно п.~5.4. \\
\hline
Среда исполнения & Компрометация среды и побочные каналы исключены; TEE не используется. & Рассмотреть отдельный профиль с TEE, пересмотрев доверие к аппаратной среде и оставшиеся угрозы согласно п.~5.3. \\
\hline
Сценарий обучения & Совместное обучение на долях; режим FL не предусмотрен. & Исследовать локальное обучение и защищённое агрегирование обновлений; отдельно определить доступ к весам и раскрытие результатов согласно п.~5.4. \\
\hline
Служебные сведения & Размерности и время не скрываются; ORAM не применяется. & Оценить допустимость утечек; при необходимости исследовать ORAM для адресов обращений. Скрытие размерностей рассматривать отдельно согласно п.~5.2. \\
\hline
\end{longtable}
}

Первым этапом после приёмки будут испытания на дополнительных данных и анализ измеренных издержек. Изменения протокола и механизмов выдачи ответа будут рассматриваться после уточнения модели угроз. TEE, FL и ORAM не будут включаться одновременно: отдельные эксперименты позволят ограничить сложность и установить вклад каждого изменения, не подменяя цель MVP максимальной эффективностью или защищённостью.

Граница блоков из подраздела~\ref{subsec:product-architecture} позволит сохранять описание модели при замене реализации защищённых операций. Режим FL потребует расширения контракта заданий и правил доступа к весам, но не отказа от разделения блоков. Каждая замена потребует повторных проверок численной совместимости и безопасности; неизменность интерфейса сама по себе их не гарантирует. Практическое внедрение дополнительно потребует организационных процедур подраздела~\ref{subsec:standardization}.

\subsection{Вывод}
\label{subsec:mvp-conclusion}

Предложен план разработки продукта из двух раздельных блоков: модели МО и заменяемого МКВ-модуля. Ожидаемый результат MVP --- демонстрация обучения и инференса логистической регрессии на данных трёх организаций без раскрытия их входов другим участникам и координатору при допущениях подраздела~\ref{subsec:mvp-threat-model}.

Выбор готовых библиотечных операций обусловлен простотой реализации, а не достижением максимальной точности, эффективности или защищённости. Выполнение критериев подраздела~\ref{subsec:mvp-implementation-plan} должно подтвердить работоспособность интеграции и соблюдение численных допусков. Демонстрация не заменит доказательство безопасности и не подтвердит готовность к промышленному внедрению.

План связывает анализ криптографических основ, классификации и применений МКВ с проверяемым экспериментом, конкретизируя цель работы и задачу подготовки будущего исследования. Разделение блоков позволит развивать продукт по направлениям подраздела~\ref{subsec:mvp-risks-roadmap} без включения криптографической логики в описание модели.

\newpage\section*{Заключение}\addcontentsline{toc}{section}{Заключение}

В работе систематизированы теоретические и прикладные аспекты МКВ. Криптографические основы обеспечения корректности и приватности изложены в разделе~\ref{sec:cryptographic-foundations}, классификация по числу участников и моделям безопасности --- в разделе~\ref{sec:protocol-classification}, применения, включая анализ зарплат, обучение нейронных сетей и управление криптовалютными ключами, --- в разделе~\ref{sec:applications}.

Сопоставление семейств протоколов приведено в таблице~\ref{tab:mpc-primitives}, ограничения вычислений, сетевого обмена и обработки дробных чисел --- в разделе~\ref{sec:limitations}. Работа Zhang et al. показывает сокращение числа умножений при сохранении точности за счёт оптимизации свёртки и Softmax на основе схемы Шамира (см. подраздел~\ref{subsec:ml-applications}); эффективность нелинейных функций остаётся нерешённой задачей.

Для будущего исследования предложен продукт из двух раздельных блоков: логистической регрессии и заменяемого МКВ-модуля. Выбраны TFEncrypted и ABY3 для трёх серверов в semi-honest модели при честном большинстве. Приоритет --- простота реализации, а не максимальная эффективность или защищённость (см. раздел~\ref{sec:future-research}).

Цель MVP --- обучение и инференс без раскрытия входов с проверкой численных допусков относительно открытой контрольной модели (см. подраздел~\ref{subsec:mvp-implementation-plan}). Реализация и измерения предстоят. Демонстрация не заменит доказательство безопасности; защита от активного противника, fairness и промышленная готовность не заявляются.

После MVP предусматриваются расширение испытаний и отдельное исследование DP, TEE, FL и ORAM (см. подраздел~\ref{subsec:mvp-risks-roadmap}). Замена МКВ-модуля при сохранении границы блоков потребует повторной проверки численной совместимости и модели угроз.

\newpage\addcontentsline{toc}{section}{Список использованных источников}\begin{thebibliography}{99}

\bibitem{01} Загартдинов Б. Н., Поляков М. В. АНАЛИЗ РЕАЛИЗАЦИИ ТЕХНОЛОГИЙ КОНФИДЕНЦИАЛЬНЫХ ВЫЧИСЛЕНИЙ //Вопросы кибербезопасности. --- 2023. --- №. 6 (58). --- С. 122-127.

\bibitem{02} Земцов Д. С. СРАВНИТЕЛЬНАЯ ОЦЕНКА СХЕМ ГОМОМОРФНОГО ШИФРОВАНИЯ НА ОСНОВЕ ЭЛЛИПТИЧЕСКОЙ КРИВОЙ ДЛЯ НОВЫХ БЕЗОПАСНЫХ МНОГОСТОРОННИХ ВЫЧИСЛЕНИЙ //Э40 Научный электронный журнал «Инновации. Наука. Образование. Отв. ред. Сафронов АИ---Тольятти:- 2024.- № 107 (октябрь).- 333 с.- URL: http://innovjourn. ru. --- 2015.

\bibitem{03} Румянцев К. Е., Хаджиева Л. К. Конфиденциальные вычисления в условиях защищенной вычислительной среды //Вестник ГГНТУ. Технические науки. --- 2022. --- Т. 18. --- №. 4 (30). --- С. 31.

\bibitem{04} Щербаков А. Ю. Конфиденциальное машинное обучение на основе двусторонних протоколов безопасных вычислений //Безопасность информационных технологий. --- 2021. --- Т. 28. --- №. 4. --- С. 39-51.

\bibitem{05} Bao H. et al. Secure multiparty computation protocol based on homomorphic encryption and its application in blockchain //Heliyon. --- 2024. --- Т. 10. --- №. 14.

\bibitem{06} Evans D., Kolesnikov V., Rosulek M. A pragmatic introduction to secure multi-party computation //Foundations and Trends® in Privacy and Security. --- 2018. --- Т. 2. --- №. 2-3. --- С. 70-246.

\bibitem{07} Liao T., Li T., Nadkarni P. A survey on secure machine learning //arXiv preprint arXiv:2505.15124. --- 2025.

\bibitem{08} Lindell Y. Secure multiparty computation //Communications of the ACM. --- 2020. --- Т. 64. --- №. 1. --- С. 86-96.

\bibitem{09} Morales D., Agudo I., Lopez J. Private set intersection: A systematic literature review //Computer Science Review. --- 2023. --- Т. 49. --- С. 100567.

\bibitem{10} Moreira P. G. Secure Multi-party Computation: A Survey : дис. --- Delft University of Technology, 2025.

\bibitem{11} Naresh V. S., Venkata Raju A., Srinivasa Rao O. Secure Multiparty Computation for Privacy-Preserving Machine Learning in Healthcare: A Comprehensive Survey //Wiley Interdisciplinary Reviews: Computational Statistics. --- 2025. --- Т. 17. --- №. 3. --- С. e70046.

\bibitem{12} Zhang S. et al. Efficient and secure multi-party computation protocol supporting deep learning //Cybersecurity. --- 2025. --- Т. 8. --- №. 1. --- С. 46.

\bibitem{13} Zhou I. et al. Secure multi-party computation for machine learning: A survey //IEEE Access. --- 2024. --- Т. 12. --- С. 53881-53899.

\bibitem{14} Yao A. C. Protocols for secure computations //Proceedings of the 23rd Annual Symposium on Foundations of Computer Science. --- IEEE, 1982. --- С. 160--164.

\bibitem{15} Yao A. C. How to generate and exchange secrets //Proceedings of the 27th Annual Symposium on Foundations of Computer Science. --- IEEE, 1986. --- С. 162--167.

\bibitem{16} Beaver D., Micali S., Rogaway P. The round complexity of secure protocols //Proceedings of the 22nd Annual ACM Symposium on Theory of Computing. --- ACM, 1990. --- С. 503--513.

\bibitem{17} Boston Women’s Workforce Council. Report 2016. --- Boston: City of Boston, 2017. --- URL: \url{https://www.boston.gov/sites/default/files/document-file-09-2017/bwwcr-2016-new-report.pdf} (дата обращения: 04.10.2026).

\bibitem{18} Cybernetica AS. Track Big Data for Governments and Education. --- 2021. --- 7 July. --- URL: \url{https://cyber.ee/resources/case-studies/track-big-data-for-governments-and-education/} (дата обращения: 04.10.2026).

\bibitem{19} MPC Alliance. Industry Leaders Launch MPC Alliance to Elevate Security and Privacy of Online Services. --- 2019. --- 13 November. --- URL: \url{https://www.mpcalliance.org/news} (дата обращения: 04.10.2026).

\bibitem{20} Mohassel P., Rindal P. ABY3: A Mixed Protocol Framework for Machine Learning //Cryptology ePrint Archive. --- 2018. --- Paper 2018/403. --- URL: \url{https://eprint.iacr.org/2018/403} (дата обращения: 05.10.2026).

\bibitem{21} TF Encrypted. TFEncrypted: репозиторий исходного кода и документация. --- URL: \url{https://github.com/tf-encrypted/tf-encrypted} (дата обращения: 05.10.2026).

\bibitem{22} Facebook Research. CrypTen: репозиторий исходного кода и документация. --- URL: \url{https://github.com/facebookresearch/CrypTen} (дата обращения: 05.10.2026).

\bibitem{23} Mohassel P., Zhang Y. SecureML: A System for Scalable Privacy-Preserving Machine Learning //Cryptology ePrint Archive. --- 2017. --- Paper 2017/396. --- URL: \url{https://eprint.iacr.org/2017/396} (дата обращения: 05.10.2026).

\bibitem{24} Cohen R., Lindell Y. Fairness Versus Guaranteed Output Delivery in Secure Multiparty Computation //Advances in Cryptology -- ASIACRYPT 2014. --- Lecture Notes in Computer Science. --- 2014. --- Т. 8874. --- С. 466--485.

\end{thebibliography}

\end{document}